\chapter{正規表現の基本}
\label{app:regex}
\index{せいきひょうげん@正規表現|textbf}

正規表現は、文字列が満たすパターンを記述する記法である。\code{grep}、\code{awk}、Python の \code{re} モジュールでは、ログの抽出、ファイル名の検査、検出器 ID の選別、時刻文字列の検証などに正規表現を用いる。ここでは、本文で触れた \code{grep} や \code{awk} の延長として、解析に必要な構文をまとめる。

以下のメタ文字の説明は、主に Python の \code{re} と \code{grep -E} の拡張正規表現を対象とする~\cite{python_re,grep_manual}。\code{grep} の既定の基本正規表現とは、\code{+}、\code{?}、括弧などの扱いが異なる。Python の \code{.} は既定では改行に一致せず、\code{\textbackslash d} は ASCII の数字以外の Unicode の十進数字にも一致する。半角数字だけが対象なら \code{[0-9]} と明示する。

\section{文字列パターンを表す正規表現}

\subsection{文字どおりの一致とメタ文字}
\index{めたもじ@メタ文字|textbf}

まず、文字どおりの検索と正規表現による検索を区別する。\code{det\_A} という文字列だけを探す場合は、各文字をそのまま並べる。

\begin{lstlisting}[numbers=none, xleftmargin=2\zw]
$ grep "det_A" events.txt
\end{lstlisting}

これに対して、\code{.} や \code{*} のような特別な意味を持つ記号を使うと、「どの文字でもよい」「0 回以上くり返す」といった条件を書けるようになる。こうした特別な意味を持つ文字をメタ文字という。

\begin{lstlisting}[numbers=none, xleftmargin=2\zw]
.
*
+
?
^
$
[]
()
|
\end{lstlisting}

たとえば \code{det\_.} は \code{det\_A} にも \code{det\_B} にも一致する。\code{.} は「任意の 1 文字」を表すからである。ピリオドそのものとの一致を指定するには、前にバックスラッシュを付け、\texttt{\textbackslash.} のようにエスケープする。\code{+} や \code{?} を文字どおりに照合する場合も同様である。

\subsection{ワイルドカードとの違い}
\index{わいるどかーど@ワイルドカード}

シェルのワイルドカードと正規表現は似て見えるが、同じものではない。\code{ls *.txt} の \code{*} は「任意の文字列」を表すが、正規表現の \code{*} は「直前の要素の 0 回以上のくり返し」を表す。したがって、正規表現で「任意の文字列」を書きたいときは \code{.*} と書く。

\begin{table}[hbtp]
\centering
\caption{ワイルドカードと正規表現の違い}
\label{tab:wildcard-vs-regex}
\begin{tabular}{lll}
\toprule
場面 & 記法 & 意味 \\
\midrule
シェル & \code{*.txt} & 任意の名前で \code{.txt} で終わるファイル \\
正規表現 & \texttt{.*\textbackslash.txt\$} & 任意の文字列で始まり \code{.txt} で終わる文字列 \\
\bottomrule
\end{tabular}
\end{table}

表~\ref{tab:wildcard-vs-regex} は、二つの記法における \code{*} の意味を比較している。\code{grep "*.txt"} では、引用符内の \code{*} をシェルではなく \code{grep} の正規表現エンジンが解釈する。正規表現の \code{*} は直前の要素の 0 回以上の反復を表すため、この位置では「任意の文字列」にならず、エラーとなる場合もある。シェルによるファイル名展開にはワイルドカード、文字列内容の照合には正規表現を用いる。

正規表現は文字列の形を検査する。日付や時刻のパターンに一致しても、2 月 30 日や 99 時などが有効な日時になるとは限らない。暦や時刻の値域は \code{datetime} などでも検証する。

\section{アンカー・文字クラス・量指定子}

\subsection{アンカー（行頭と行末）}
\index{あんかー@アンカー（正規表現）|textbf}

\code{\textasciicircum} は行頭、\code{\$} は行末を表す。したがって、文字列の一部ではなく「行全体がその形である」ことを確かめたいときに使う。

\begin{lstlisting}[numbers=none, xleftmargin=2\zw]
$ grep '^ERROR' run.log
$ grep 'txt$' filelist.txt
$ grep '^det_A$' detectors.txt
\end{lstlisting}

1 行目は \code{ERROR} で始まる行、2 行目は \code{txt} で終わる行、3 行目は行全体がちょうど \code{det\_A} である行だけを取る。特に 3 行目のように両端をアンカーで挟むと、\code{det\_AB} や \code{xdet\_A} が紛れ込むのを防げる。

\subsection{文字クラス \code{[...]} と代表的な省略記法}
\index{もじくらす@文字クラス|textbf}

\code{[abc]} は \code{a}、\code{b}、\code{c} のいずれか 1 文字、\code{[A-Z]} は英大文字 1 文字を表す。検出器名やセンサー番号のように、限られた候補から 1 文字を選ぶパターンに用いる。

\begin{lstlisting}[numbers=none, xleftmargin=2\zw]
^det_[AB]$
^run_[0-9][0-9][0-9]$
^[A-Za-z_][A-Za-z0-9_]*$
\end{lstlisting}

1 行目は \code{det\_A} または \code{det\_B}、2 行目は \code{run\_001} のような 3 桁番号付きの名前、3 行目は Python 風の識別子に近い形を表す。範囲指定の \code{[0-9]} は、数字を含むパターンで頻出する。

Python の \code{re} では、数字 1 文字を \code{\textbackslash d}、空白 1 文字を \code{\textbackslash s} と書ける。

\begin{lstlisting}[language=Python]
import re

re.fullmatch(r"det_[AB]", "det_A")
re.fullmatch(r"\d\d:\d\d:\d\d", "12:34:56")
re.search(r"\s+", "det_A   15.2")
\end{lstlisting}

ただし、標準的な \code{grep} の基本・拡張正規表現と \code{awk} では、\code{\textbackslash d} や \code{\textbackslash s} を Python と同じ省略記法として使用できない。コマンドラインでは \code{[0-9]} や \code{[[:space:]]} を用いる。

\subsection{量指定子 \code{*}, \code{+}, \code{?}}
\index{りょうしていし@量指定子|textbf}

量指定子は、直前の要素が何回くり返されるかを表す。\code{*} は 0 回以上、\code{+} は 1 回以上、\code{?} は 0 回または 1 回である。

\begin{lstlisting}[numbers=none, xleftmargin=2\zw]
ab*c
ab+c
colou?r
\end{lstlisting}

\code{ab*c} は \code{ac}、\code{abc}、\code{abbbc} に一致する。\code{ab+c} は少なくとも 1 個の \code{b} が必要なので \code{ac} には一致しない。\code{colou?r} は \code{color} と \code{colour} の両方に一致する。

量指定子を使うと、1 個以上の空白、任意指定の符号、任意指定の小数部など、要素の出現回数をパターンに記述できる。

\begin{lstlisting}[language=Python]
pattern = re.compile(r"^[+-]?\d+(\.\d+)?$")

for text in ["12", "-3.5", "+0.25", "1.2.3"]:
    print(text, bool(pattern.fullmatch(text)))
\end{lstlisting}

このパターンは、符号付き整数や小数を表し、\code{1.2.3} のような壊れた文字列は弾く。

\section{run 名・時刻・数値ログの抽出}

\subsection{ログからの対象行抽出}

実験ログやジョブログでは、異常行や特定の run だけを抽出する場合がある。アンカーと文字クラスを組み合わせると、行全体に対する条件を具体的に記述できる。

\begin{lstlisting}[numbers=none, xleftmargin=2\zw]
$ grep '^ERROR' run.log
$ grep '^2026-04-04 .* det_[AB] ' run.log
$ grep -E '^(WARN|ERROR)' run.log
\end{lstlisting}

3 行目の \code{grep -E} は、拡張正規表現として \code{|} による選択を解釈する。拡張正規表現では、括弧、\code{+}、\code{?}、\code{|} をバックスラッシュなしで記述できる。

\subsection{検出器 ID・センサー名の選別}

検出器名が \code{det\_A}, \code{det\_B}, \code{det\_C} のように並んでいるなら、文字クラスや選択を使って対象を絞れる。

\begin{lstlisting}[numbers=none, xleftmargin=2\zw]
$ grep '^det_[AB]$' detector_ids.txt
$ awk '$1 ~ /^det_[AB]$/ {print $1, $3}' summary.txt
\end{lstlisting}

\code{awk} の \code{\$1 ~ /.../} は、第 1 列が正規表現に一致するかを判定する。列ごとに異なる照合条件や出力列を指定する場合は、列を直接参照できる \code{awk} を用いる。

\subsection{Python による抽出と変換}

正規表現の括弧で囲んだ部分は、照合結果から個別に取り出せる。Python では、\code{re.search} や \code{re.fullmatch} の返り値に対して \code{group()} を呼び出す。

\begin{lstlisting}[language=Python]
import re

line = "2026-04-04 12:34:56 det_A signal=15.2"
match = re.search(r"signal=(\d+\.\d+)", line)

if match:
    signal = float(match.group(1))
    print(signal)
\end{lstlisting}

括弧 \code{(...)} で囲んだ部分はキャプチャされ、\code{group(1)}、\code{group(2)} のように取り出せる。ログ 1 行から run 番号、検出器名、信号値を個別に抽出する場合に用いる。

\section{Python の \code{re} と \code{grep}/\code{awk} の使い分け}

\subsection{Python の raw string によるパターン記述}
\index{raw string@raw string|textbf}
\index{re module@\code{re} モジュール|textbf}

Python の通常の文字列リテラルでは、バックスラッシュが Python と正規表現の双方で特別な意味を持つ。正規表現は、Python 側でバックスラッシュを解釈しない \code{r"..."} 形式の raw string で記述する。

\begin{lstlisting}[language=Python]
import re

pattern = re.compile(r"^det_[A-Z]\d+$")
print(bool(pattern.fullmatch("det_A12")))
\end{lstlisting}

通常の Python 文字列では、\texttt{"\textbackslash\textbackslash d+"} のようにバックスラッシュを二重に記述する必要がある。正規表現のパターンは、原則として raw string で記述する。

\subsection{\code{search}, \code{match}, \code{fullmatch} の使い分け}
\index{re module@\code{re} モジュール!search@\code{search}}
\index{re module@\code{re} モジュール!match@\code{match}}
\index{re module@\code{re} モジュール!fullmatch@\code{fullmatch}}

\code{re.search} は文字列中の任意の位置、\code{re.match} は先頭、\code{re.fullmatch} は文字列全体に対して一致を調べる。ファイル名や識別子の全体を検証する場合は、接頭辞または接尾辞に余分な文字があれば不一致となる \code{fullmatch} を用いる。

\begin{lstlisting}[language=Python]
import re

print(re.search(r"\d+", "run_42").group())
print(bool(re.match(r"run", "run_42")))
print(bool(re.fullmatch(r"run_\d+", "run_42")))
\end{lstlisting}

\subsection{コマンドラインで十分な場面と、Python に移すべき場面}

条件に合う行の表示、件数の確認、特定列の抽出は、\code{grep} または \code{awk} だけで記述できる。一致部分を数値へ変換して計算する場合、複数の条件分岐を組み合わせる場合、または抽出結果から DataFrame を作る場合は、変換、分岐、テストを同じコードに記述できる Python を用いる。

\begin{itemize}
\item \code{grep}: 行全体の抽出や件数確認をすばやく行いたいとき
\item \code{awk}: 列条件と正規表現を組み合わせたいとき
\item Python \code{re}: 抽出した結果をそのまま解析処理へつなげたいとき
\end{itemize}

正規表現は必要な条件を小さく組み合わせて作る。受理する文字列と拒否する文字列のテストを用意し、変更時に再確認する。
